\documentclass[10pt,a4paper]{article}

% ==============================================================================
% TEMPLATE EDITORIAL — BOLETIM FLASH (JORNAL ECONÔMICO DO TOCANTINS)
% Estilo: Briefing Executivo de Jornal Econômico (2 Páginas Balanceadas)
% Paleta FLASH: Azul-Marinho (#0F2537), Azul Secundário (#244B6B),
%               Dourado Institucional (#C8942B) e Fundo Neutro (#F4F6F9)
% ==============================================================================
\usepackage[utf8]{inputenc}
\usepackage[T1]{fontenc}
\usepackage[brazil]{babel}
\usepackage{helvet}
\renewcommand{\familydefault}{\sfdefault}

\usepackage[top=1.35cm, bottom=1.45cm, left=1.40cm, right=1.40cm, headheight=14pt, headsep=8pt, footskip=18pt]{geometry}
\usepackage{xcolor}
\usepackage{graphicx}
\usepackage{booktabs}
\usepackage{tabularx}
\usepackage{longtable}
\usepackage{array}
\usepackage{colortbl}
\usepackage{tikz}
\usepackage{fancyhdr}
\usepackage{lastpage}
\usepackage{enumitem}
\usepackage{microtype}
\usepackage{setspace}
\usepackage{multicol}
\usepackage{amssymb}
\usepackage{hyperref}

% Cores institucionais do Boletim FLASH
\definecolor{corprimaria}{HTML}{0F2537}
\definecolor{corsecundaria}{HTML}{244B6B}
\definecolor{corfundo}{HTML}{F8FAFC}
\definecolor{douradoto}{HTML}{C8942B}
\definecolor{coracento}{HTML}{C8942B}
\definecolor{cinzatexto}{HTML}{1E293B}
\definecolor{cinzasub}{HTML}{475569}
\definecolor{cinzaborda}{HTML}{CBD5E1}
\definecolor{verdealta}{HTML}{15803D}
\definecolor{vermelhobaixa}{HTML}{B91C1C}

\hypersetup{
  colorlinks=true,
  linkcolor=corprimaria,
  urlcolor=corsecundaria,
  pdftitle={Mercado de trabalho formal no Tocantins mantém geração líquida positiva puxada por Serviços e Construção},
  pdfauthor={Jornal Economico do Tocantins}
}

\setlength{\parindent}{0pt}
\setlength{\parskip}{3.5pt}
\setlength{\columnsep}{0.55cm}
\setlength{\columnseprule}{0.4pt}
\def\columnseprulecolor{\color{cinzaborda}}
\renewcommand{\arraystretch}{1.15}

% Cabeçalho e Rodapé de Jornal Econômico (Sem colisão central)
\pagestyle{fancy}
\fancyhf{}
\renewcommand{\headrulewidth}{0pt}
\renewcommand{\footrulewidth}{0.5pt}
\renewcommand{\footrule}{\hbox to\headwidth{\color{douradoto}\leaders\hrule height \footrulewidth\hfill}}

\fancyhead[L]{\ifnum\value{page}>1 \scriptsize\bfseries\color{corprimaria}JORNAL ECONÔMICO DO TOCANTINS \enspace{\color{douradoto}$\vert$}\enspace {\mdseries\color{cinzasub}BOLETIM FLASH DE CONJUNTURA $\cdot$ Outubro/2026}\fi}
\fancyhead[R]{\ifnum\value{page}>1 \scriptsize\color{cinzasub}\#03/2026\fi}

\fancyfoot[L]{\scriptsize\color{cinzasub}\textbf{\color{corprimaria}Jornal Econômico do Tocantins} $\cdot$ Mercado de Trabalho: Novo CAGED, RAIS e PNAD Contínua (TO) $\cdot$ Publicado em 29/10/2026}
\fancyfoot[R]{\scriptsize\bfseries\color{corprimaria}Pág. \thepage\ de \pageref*{LastPage}}

% Título de Seção Editorial Limpo (Fio fino inferior, sem bloco escuro pesado)
\newcommand{\secaoboletim}[1]{%
  \par\vspace{0.18cm}%
  \noindent{\fontsize{9.0}{10.8}\selectfont\bfseries\color{corprimaria}\MakeUppercase{#1}}%
  \par\vspace{2pt}%
  {\color{douradoto}\hrule height 1.0pt}%
  \vspace{0.08cm}%
}

\newcommand{\subsecaoboletim}[1]{%
  \par\vspace{0.10cm}%
  \noindent{\small\bfseries\color{corprimaria}#1}%
  \par\vspace{0.03cm}%
  {\color{cinzaborda}\hrule height 0.5pt}%
  \vspace{0.06cm}%
}

\begin{document}
\color{cinzatexto}

% ==============================================================================
% MASTHEAD EDITORIAL CLARO ESTILO PERIÓDICO FINANCEIRO (ZERO BLOCOS PESADOS)
% ==============================================================================
\noindent{\color{corprimaria}\hrule height 2.2pt}%
\par\vspace{2pt}%
{\color{douradoto}\hrule height 0.8pt}%
\par\vspace{6pt}%
\noindent\hbox to \linewidth{%
  {\scriptsize\bfseries\color{douradoto}\MakeUppercase{BOLETIM FLASH DE CONJUNTURA} \enspace$\cdot$\enspace \#03/2026}%
  \hfill
  {\scriptsize\color{cinzasub}\textbf{\color{corprimaria}JORNAL ECONÔMICO DO TOCANTINS} \enspace$\vert$\enspace Ref.: \textbf{Outubro/2026} \enspace$\vert$\enspace 29/10/2026}%
}%
\par\vspace{5pt}%
{\fontsize{14.5}{17.4}\selectfont\bfseries\color{corprimaria}Mercado de trabalho formal no Tocantins mantém geração líquida positiva puxada por Serviços e Construção\par}%
\vspace{3pt}%
{\footnotesize\itshape\color{cinzasub}Dinamismo nos polos regionais de Palmas, Araguaína e Gurupi sustenta baixa taxa de desocupação estadual\par}%
\par\vspace{6pt}%
{\color{cinzaborda}\hrule height 0.6pt}%

\vspace{0.14cm}

% ==============================================================================
% 1. FAIXA DE INDICADORES-CHAVE DO PAINEL (4 CARDS KPI EM LINHA ÚNICA)
% ==============================================================================
\noindent\makebox[\textwidth][s]{%
\begin{tikzpicture}[baseline=(box.base)]
  \node[fill=white, draw=cinzaborda, line width=0.35pt,
        inner xsep=6pt, inner ysep=4pt, text width=0.206\textwidth, minimum height=1.14cm, align=left] (box) {%
    {\fontsize{6.0}{7.2}\selectfont\bfseries\color{cinzasub}ESTOQUE RAIS (2025)}\par\vspace{1.2pt}%
    {\fontsize{9.2}{10.8}\selectfont\bfseries\color{corprimaria}437.580}\par\vspace{1pt}%
    {\fontsize{5.8}{6.8}\selectfont\color{cinzasub}Estab. Ativos: 33.198}%
  };
  \fill[coracento] (box.north west) rectangle ([yshift=-1.6pt]box.north east);
\end{tikzpicture}\hfill%
\begin{tikzpicture}[baseline=(box.base)]
  \node[fill=white, draw=cinzaborda, line width=0.35pt,
        inner xsep=6pt, inner ysep=4pt, text width=0.206\textwidth, minimum height=1.14cm, align=left] (box) {%
    {\fontsize{6.0}{7.2}\selectfont\bfseries\color{cinzasub}SALDO CAGED (09/26)}\par\vspace{1.2pt}%
    {\fontsize{9.2}{10.8}\selectfont\bfseries\color{corprimaria}+716}\par\vspace{1pt}%
    {\fontsize{5.8}{6.8}\selectfont\color{cinzasub}Adm: 11.263 | Desl: 10.547}%
  };
  \fill[coracento] (box.north west) rectangle ([yshift=-1.6pt]box.north east);
\end{tikzpicture}\hfill%
\begin{tikzpicture}[baseline=(box.base)]
  \node[fill=white, draw=cinzaborda, line width=0.35pt,
        inner xsep=6pt, inner ysep=4pt, text width=0.206\textwidth, minimum height=1.14cm, align=left] (box) {%
    {\fontsize{6.0}{7.2}\selectfont\bfseries\color{cinzasub}DESOCUPAÇÃO (PNAD)}\par\vspace{1.2pt}%
    {\fontsize{9.2}{10.8}\selectfont\bfseries\color{corprimaria}4,1\%}\par\vspace{1pt}%
    {\fontsize{5.8}{6.8}\selectfont\color{cinzasub}IBGE Tabela 4099 (27 UFs)}%
  };
  \fill[coracento] (box.north west) rectangle ([yshift=-1.6pt]box.north east);
\end{tikzpicture}\hfill%
\begin{tikzpicture}[baseline=(box.base)]
  \node[fill=white, draw=cinzaborda, line width=0.35pt,
        inner xsep=6pt, inner ysep=4pt, text width=0.206\textwidth, minimum height=1.14cm, align=left] (box) {%
    {\fontsize{6.0}{7.2}\selectfont\bfseries\color{cinzasub}REMUNERAÇÃO RAIS / PNAD}\par\vspace{1.2pt}%
    {\fontsize{9.2}{10.8}\selectfont\bfseries\color{corprimaria}R\$ 4.022,93}\par\vspace{1pt}%
    {\fontsize{5.8}{6.8}\selectfont\color{cinzasub}PNAD Habitual: R\$ 3.351}%
  };
  \fill[coracento] (box.north west) rectangle ([yshift=-1.6pt]box.north east);
\end{tikzpicture}%
}


\vspace{0.10cm}

% ==============================================================================
% 2. ANÁLISE PRELIMINAR AUTOMATIZADA DO SCRIPT (RASTREAMENTO DE VARIAÇÕES)
% ==============================================================================
\par\vspace{0.05cm}
\noindent{\color{cinzaborda}\hrule height 0.5pt}\par\vspace{4pt}
{\fontsize{7.8}{9.4}\selectfont\bfseries\color{corprimaria}\MakeUppercase{Síntese dos Indicadores no Período}}\par\vspace{2pt}
{\scriptsize\color{cinzatexto}
\begin{itemize}[leftmargin=10pt, itemsep=1.2pt, topsep=1pt, parsep=0pt]
   \item \textbf{\color{corprimaria}Saldo Mensal do Novo CAGED (09/26):} Saldo Mensal do Novo CAGED (09/26) sofreu retração acentuada de -42,40\% (queda drástica que requer atenção conjuntural). (patamar atual: +716 vagas)
   \item \textbf{\color{corprimaria}Estoque de Vínculos Ativos (RAIS):} Estoque de Vínculos Ativos (RAIS) apresentou crescimento robusto de +5,40\% no período. (patamar atual: 437.580)
   \item \textbf{\color{corprimaria}Taxa de Desocupação PNAD (Tocantins):} Taxa de Desocupação PNAD (Tocantins) apresentou acomodação negativa de -0,30\%. (patamar atual: 4,1\%)
   \item \textbf{\color{corprimaria}Rendimento Médio Real Habitual (PNAD):} Rendimento Médio Real Habitual (PNAD) avançou +4,20\% em trajetória de expansão gradual. (patamar atual: R\$ 3.351)
\end{itemize}
}
\noindent{\color{cinzaborda}\hrule height 0.5pt}\par\vspace{0.05cm}


\vspace{0.08cm}

% ==============================================================================
% 3. ANÁLISE CONJUNTURAL DA EQUIPE E PERSPECTIVAS (2 COLUNAS DE JORNAL)
% ==============================================================================
\secaoboletim{1. Análise Conjuntural da Equipe e Cenários}

\begin{multicols}{2}
\begin{spacing}{1.06}
{\small O comportamento recente do mercado de trabalho formal no Tocantins reflete um ciclo virtuoso de transmissão de renda do setor primário e das obras de infraestrutura para as atividades urbanas. A liderança do setor de Serviços e da Construção Civil na abertura líquida de postos com carteira assinada evidencia que o crescimento econômico estadual não está restrito ao campo, mas transborda para o consumo das famílias e para o mercado imobiliário nas principais cidades.

\medskip
A análise regionalizada dos municípios confirma o papel de Palmas, Araguaína e Gurupi como grandes absorvedores de mão de obra formal, seguidos por Porto Nacional e Paraíso do Tocantins, que se beneficiam diretamente de sua posição estratégica como entroncamentos logísticos e industriais no eixo da BR-153.

\medskip
Vale destacar, sob a ótica qualitativa, que a manutenção da taxa de desocupação em patamar historicamente reduzido (5,4\%) começa a gerar pressões pontuais por mão de obra qualificada em segmentos técnicos, exigindo maior articulação entre o setor produtivo e as instituições de qualificação profissional no Estado.}

\subsecaoboletim{Perspectivas e Cenários}
{\small Para a reta final do ano, a expectativa da equipe é de manutenção de saldos positivos em outubro e novembro, impulsionados pelas contratações temporárias de fim de ano no comércio varejista e nos serviços de alimentação e logística, antes do ajuste sazonal típico do mês de dezembro.}

\end{spacing}
\end{multicols}

% ==============================================================================
% 4. PANORAMA ESTATÍSTICO, GRÁFICO E TABELAS EXECUTIVAS DO PERÍODO
% ==============================================================================
\secaoboletim{2. Panorama Estatístico, Gráfico e Tabelas Executivas do Período}

\begin{center}\includegraphics[width=0.95\textwidth]{figuras/2026_10_flash_caged_trab_caged.pdf}\end{center}
\par\vspace{0.08cm}
\noindent\begin{minipage}[t]{0.482\textwidth}
\par\vspace{0.08cm}
\noindent\begin{minipage}{\linewidth}
  \noindent{\fontsize{7.8}{9.5}\selectfont\bfseries\color{corprimaria}\rule{0.10cm}{0.20cm}\hspace{4pt}Movimentação Setorial do Novo CAGED no Tocantins (Setembro/2026)}\par\vspace{2pt}
  {\scriptsize\setlength{\tabcolsep}{4.5pt}%
\noindent\resizebox{\linewidth}{!}{%
\begin{tabular}{lrrrrr}
\toprule
\rowcolor{corprimaria} \textbf{\color{white}Grande Grupamento Setorial} & \textbf{\color{white}Admissões} & \textbf{\color{white}Desligamentos} & \textbf{\color{white}Saldo Mês} & \textbf{\color{white}Saldo Acum. Ano} & \textbf{\color{white}Salário Médio Adm.} \\
\midrule
Serviços & 4.735 & 4.145 & +590 & +2.634 & R\$ 2.055,59 \\
\rowcolor{corfundo} Construção & 1.212 & 1.001 & +211 & +2.229 & R\$ 2.531,37 \\
Comércio & 3.033 & 2.912 & +121 & +329 & R\$ 1.900,07 \\
\rowcolor{corfundo} Agropecuária & 1.179 & 1.366 & -187 & -851 & R\$ 2.438,64 \\
Indústria & 1.104 & 1.123 & -19 & +788 & R\$ 2.172,38 \\
\bottomrule
\end{tabular}%
}%
}
  \par\vspace{1.5pt}{\fontsize{6.0}{7.2}\selectfont\color{cinzasub}\textit{Fonte: MTE — Microdados do Novo CAGED (MOV + FOR - EXC).}}
\end{minipage}
\par\vspace{0.08cm}
\end{minipage}\hfill
{\color{cinzaborda}\vrule width 0.4pt}\hfill
\begin{minipage}[t]{0.482\textwidth}
\par\vspace{0.08cm}
\noindent\begin{minipage}{\linewidth}
  \noindent{\fontsize{7.8}{9.5}\selectfont\bfseries\color{corprimaria}\rule{0.10cm}{0.20cm}\hspace{4pt}Top 5 Municípios do Tocantins em Estoque Formal (RAIS) e Saldo Mensal (Novo CAGED)}\par\vspace{2pt}
  {\scriptsize\setlength{\tabcolsep}{4.0pt}%
\noindent\resizebox{\linewidth}{!}{%
\begin{tabular}{clrr}
\toprule
\rowcolor{corprimaria} \textbf{\color{white}Pos.} & \textbf{\color{white}Município (Top 5 TO)} & \textbf{\color{white}Estoque (RAIS) 2025} & \textbf{\color{white}Saldo (Novo CAGED) 09/26} \\
\midrule
1\textordmasculine{} & Palmas & 220.673 & +399 \\
\rowcolor{corfundo} 2\textordmasculine{} & Araguaína & 47.010 & +180 \\
3\textordmasculine{} & Gurupi & 20.492 & +89 \\
\rowcolor{corfundo} 4\textordmasculine{} & Porto Nacional & 14.172 & +76 \\
5\textordmasculine{} & Paraíso do Tocantins & 12.288 & +27 \\
\bottomrule
\end{tabular}%
}%
}
  \par\vspace{1.5pt}{\fontsize{6.0}{7.2}\selectfont\color{cinzasub}\textit{Fonte: MTE — Microdados da RAIS (vínculos ativos em 31/12 sem abandonados) e Novo CAGED.}}
\end{minipage}
\par\vspace{0.08cm}
\end{minipage}
\par\vspace{0.08cm}

\par\vspace{0.08cm}
\noindent\begin{minipage}{\linewidth}
  \noindent{\fontsize{7.8}{9.5}\selectfont\bfseries\color{corprimaria}\rule{0.10cm}{0.20cm}\hspace{4pt}PNAD Contínua (IBGE - Tabela 4099): Ranking Completo da Taxa de Desocupação nas 27 UFs e Brasil (1\textordmasculine{} = Menor Desemprego)}\par\vspace{2pt}
  {\scriptsize\setlength{\tabcolsep}{4.5pt}%
\noindent\resizebox{\linewidth}{!}{%
\begin{tabular}{clrrclrr}
\toprule
\rowcolor{corprimaria} \textbf{\color{white}Pos.} & \textbf{\color{white}Unidade da Federação} & \textbf{\color{white}Taxa} & \textbf{\color{white}Var.} & \textbf{\color{white}Pos.} & \textbf{\color{white}Unidade da Federação} & \textbf{\color{white}Taxa} & \textbf{\color{white}Var.} \\
\midrule
1\textordmasculine{} & Santa Catarina & 2,1\% & -0.6 p.p. (trim) & 15\textordmasculine{} & Maranhão & 6,0\% & -0.9 p.p. (trim) \\
\rowcolor{corfundo} 2\textordmasculine{} & Mato Grosso & 2,2\% & -0.9 p.p. (trim) & 16\textordmasculine{} & Pará & 6,2\% & -0.8 p.p. (trim) \\
3\textordmasculine{} & Espírito Santo & 2,3\% & -0.9 p.p. (trim) & 17\textordmasculine{} & Distrito Federal & 6,5\% & -0.6 p.p. (trim) \\
\rowcolor{corfundo} 4\textordmasculine{} & Rondônia & 2,6\% & -1.1 p.p. (trim) & 18\textordmasculine{} & Acre & 6,6\% & -1.6 p.p. (trim) \\
5\textordmasculine{} & Mato Grosso do Sul & 2,7\% & -1.1 p.p. (trim) & 19\textordmasculine{} & Ceará & 6,6\% & -0.7 p.p. (trim) \\
\rowcolor{corfundo} 6\textordmasculine{} & Paraná & 3,1\% & -0.4 p.p. (trim) & 20\textordmasculine{} & Rio de Janeiro & 7,1\% & -0.2 p.p. (trim) \\
7\textordmasculine{} & Minas Gerais & 3,8\% & -1.2 p.p. (trim) & 21\textordmasculine{} & Amazonas & 7,5\% & -0.8 p.p. (trim) \\
\rowcolor{corfundo} 8\textordmasculine{} & Goiás & 4,0\% & -1.1 p.p. (trim) & 22\textordmasculine{} & Alagoas & 7,9\% & -1.3 p.p. (trim) \\
9\textordmasculine{} & * Tocantins (TO) & 4,1\% & -1.5 p.p. (trim) & 23\textordmasculine{} & Sergipe & 7,9\% & -0.7 p.p. (trim) \\
\rowcolor{corfundo} 10\textordmasculine{} & Rio Grande do Sul & 4,2\% & +0.2 p.p. (trim) & 24\textordmasculine{} & Piauí & 8,3\% & -0.6 p.p. (trim) \\
11\textordmasculine{} & Roraima & 5,0\% & -0.7 p.p. (trim) & 25\textordmasculine{} & Pernambuco & 8,3\% & -0.9 p.p. (trim) \\
\rowcolor{corfundo} 12\textordmasculine{} & Paraíba & 5,4\% & -1.6 p.p. (trim) & 26\textordmasculine{} & Bahia & 9,1\% & -0.1 p.p. (trim) \\
13\textordmasculine{} & São Paulo & 5,4\% & -0.6 p.p. (trim) & 27\textordmasculine{} & Amapá & 9,8\% & -0.2 p.p. (trim) \\
\rowcolor{corfundo} 14\textordmasculine{} & Rio Grande do Norte & 5,6\% & -2.0 p.p. (trim) & BR & Brasil & 5,4\% & -0.7 p.p. (trim) \\
\bottomrule
\end{tabular}%
}%
}
  \par\vspace{1.5pt}{\fontsize{6.0}{7.2}\selectfont\color{cinzasub}\textit{Fonte: IBGE — PNAD Contínua Trimestral (Tabela 4099). Nota: Variação em pontos percentuais (p.p.).}}
\end{minipage}
\par\vspace{0.08cm}



\vfill

% ==============================================================================
% EXPEDIENTE EDITORIAL E NOTA METODOLÓGICA
% ==============================================================================
\noindent\begin{tikzpicture}
  \node[
    fill=corfundo,
    draw=cinzaborda,
    line width=0.5pt,
    rounded corners=2pt,
    inner sep=6pt,
    text width=\textwidth-13pt,
    font=\scriptsize\color{cinzasub}
  ] {
    \textbf{\color{corprimaria}Nota Metodológica e Fontes Primárias:} Consolida os microdados mensais do Novo CAGED (MOV + FOR - EXC), o estoque anual de vínculos ativos da RAIS em 31/12 (excluídos vínculos abandonados) para os Top 5 municípios e o ranking completo da taxa de desocupação nas 27 UFs e Brasil (Tabela 4099 do IBGE).\par\vspace{2pt}
    \textbf{\color{corprimaria}Expediente Editorial:} Tarcisio Iago | Economista — Jornal Econômico do Tocantins \enspace$\vert$\enspace \textbf{\color{corprimaria}Publicação:} Núcleo de Conjuntura Econômica — \textit{Jornal Econômico do Tocantins}.
  };
\end{tikzpicture}

\end{document}
